% ---------------------------------------------------------------------------
\section{Trabalhos Relacionados}
\label{sec:trabalhos_relacionados}
% ---------------------------------------------------------------------------

A aplicação de Modelos de Linguagem de Grande Porte (LLMs) à Geração Procedural de Conteúdo (PCG) constitui um campo de investigação recente, emergente e em rápida consolidação na literatura científica de jogos digitais e computação gráfica entre os anos de 2023 e 2026 \cite{Gallotta_2024, maleki2024procedural}. O atrativo central desse paradigma reside na capacidade demonstrada pelos transformadores de interpretar instruções abertas em linguagem natural, compreender relações semânticas complexas e sintetizar estruturas conceituais coerentes a partir de contextos mínimos de entrada \cite{todd2023level, Gallotta_2024}. Contudo, a transposição direta dessas competências conexionistas para a produção de ambientes tridimensionais (3D) navegáveis e funcionais esbarra em limitações estruturais severas, em especial na incapacidade dos modelos autorregressivos de realizar raciocínio métrico, espacial e geométrico contínuo com precisão exata \cite{xu2025constraint, feng2026repurposing3dgenerativemodel}.

Para situar a pesquisa no estado da arte contemporâneo e fundamentar as decisões de projeto adotadas neste trabalho, esta seção organiza as propostas correlatas em quatro famílias arquiteturais fundamentais, acrescidas de uma análise de diagnósticos empíricos e da percepção da indústria de desenvolvimento:
\begin{enumerate}[label=(\roman*)]
    \item \textit{Geração Simbólica e Síntese de Ambientes 2D}: abordagens que exploram a capacidade \textit{in-context} de LLMs sobre representações matriciais discretas (grades ASCII, \textit{tilemaps} e triplas semânticas);
    \item \textit{Síntese de Cenas 3D via Código Intermediário e Renderização Offline}: sistemas que empregam LLMs como geradores de scripts imperativos (como código Python para Blender) combinados a modelos multimodais de visão computacional em ciclos de retroalimentação;
    \item \textit{Planejamento de Layouts 3D Guiado por Otimizadores e Restrições}: propostas que utilizam LLMs para mapear intenções textuais em linguagens de restrições formais ou modelos desacoplados, delegando a geometria a solucionadores numéricos ou heurísticos de reparo;
    \item \textit{Integração com Motores Gráficos e Arquiteturas Multiagente}: pesquisas que acoplam LLMs diretamente a motores de jogos comerciais (\textit{engines}) ou organizam redes de agentes cooperativos em funções de geração e crítica textual.
\end{enumerate}

Cada trabalho correlato é submetido a uma análise dialética tripartite composta por: (1) \textbf{Mecanismo}, descrevendo o arcabouço técnico e o fluxo operacional empregado; (2) \textbf{Mérito}, destacando as contribuições científicas e avanços perante métodos anteriores; e (3) \textbf{Ponto de Ruptura}, dissecando analiticamente o gargalo metodológico ou a falha técnica impeditiva perante o problema de síntese de fases tridimensionais completas, jogáveis e com solvabilidade garantida.

% ---------------------------------------------------------------------------
\subsection{Geração Simbólica e Síntese de Ambientes 2D via Modelos de Linguagem}
\label{sec:relacionados_2d}
% ---------------------------------------------------------------------------

As primeiras investigações sistemáticas sobre o uso de LLMs para a síntese estrutural de fases concentraram-se no domínio bidimensional discreto, no qual a geometria é representada de forma matricial por caracteres ou índices alfanuméricos.

No trabalho pioneiro de Todd et al. \cite{todd2023level}, intitulado \textit{Level Generation Through Large Language Models}, os autores avaliaram a capacidade de modelos de linguagem de grande porte pré-treinados (como GPT-3) gerarem níveis completos para o jogo de quebra-cabeça clássico \textit{Sokoban}. 
\begin{itemize}
    \item \textbf{Mecanismo:} Os autores converteram as fases em representações textuais em formato de grade ASCII, em que cada caractere codifica um elemento específico do jogo (paredes, chão transitável, caixas, alvos e posição do jogador). Por meio do aprendizado em contexto (\textit{in-context learning}), pares de exemplos de entrada e saída foram fornecidos no \textit{prompt} para instruir o modelo a completar e sintetizar novas matrizes caractere a caractere de maneira autorregressiva direta.
    \item \textbf{Mérito:} O estudo comprovou experimentalmente, de maneira pioneira, que os LLMs detêm a capacidade latente de assimilar e reproduzir a gramática superficial e as regularidades estilísticas locais de níveis de jogos sem a necessidade de ajuste fino supervisionado (\textit{fine-tuning}) em bases proprietárias.
    \item \textbf{Ponto de Ruptura:} A geração puramente autorregressiva opera sob a ausência de invariantes geométricos e espaciais globais. À medida que o tamanho da grade aumenta, a deriva estocástica token a token acumula inconsistências métricas, resultando em barreiras incompletas, alvos inacessíveis e uma taxa expressiva de fases insolúveis. Além disso, a abordagem é estritamente confinada a espaços 2D discretos regulares, sendo estruturalmente incapaz de extrapolar para coordenadas métricas contínuas em $\mathbb{R}^3$, volumes de colisão ou malhas poligonais de navegação.
\end{itemize}

Ampliando o escopo para a geração procedimental integrada entre narrativa e espaço, Nasir, James e Togelius \cite{nasir2024word2worldgeneratingstoriesworlds} apresentaram o sistema \textit{Word2World: Generating Stories and Worlds through Large Language Models}.
\begin{itemize}
    \item \textbf{Mecanismo:} O sistema emprega um pipeline modular no qual LLMs sintetizam recursivamente narrativas descritivas e convertem esses enredos em mundos bidimensionais subdivididos em blocos (\textit{tile-based}), interligados por pontos de interesse narrativo. Para atestar a viabilidade básica das áreas geradas, o sistema acopla um algoritmo de busca heurística $A^*$ tradicional que verifica a existência de caminhos livres sobre a matriz discreta de tiles ortogonais.
    \item \textbf{Mérito:} A pesquisa estabelece um avanço relevante na coesão entre enredo e mapa de jogo, demonstrando a utilidade de submeter a saída textual do LLM a um validador clássico de conectividade física.
    \item \textbf{Ponto de Ruptura:} O escopo espacial permanece estritamente restrito a \textit{tilemaps} 2D estáticos e ortogonais. Não há qualquer tratamento de transformações afins tridimensionais (translação, rotação de guinada e escala volumétrica), relevo vertical ou detecção de colisão contínua. Ademais, a validação por $A^*$ restringe-se à conectividade geométrica estática de caminhos livres, ignorando a dinâmica temporal de dependências mecânicas (como sistemas de chaves e portas trancadas, onde o estado do agente altera dinamicamente a transitabilidade de barreiras no espaço).
\end{itemize}

Em outra vertente 2D, Chen e Jhala \cite{chen2025narrative} propuseram o arcabouço \textit{Narrative-to-Scene Generation: An LLM-Driven Pipeline for 2D Game Environments}.
\begin{itemize}
    \item \textbf{Mecanismo:} O pipeline traduz histórias e descrições textuais livres em um conjunto intermediário de triplas semânticas relacionais (da forma $\langle$\textit{entidade}, \textit{relação}, \textit{entidade}$\rangle$, como ``árvore à esquerda da cabana''). Em seguida, essas relações espaciais qualitativas são traduzidas em restrições de um Problema de Satisfação de Restrições (CSP) bidimensional, resolvido algoritmicamente para dispor os elementos visuais sobre a cena.
    \item \textbf{Mérito:} O trabalho mitiga a fragilidade numérica do LLM ao abstrair coordenadas absolutas em relações qualitativas estruturadas, delegando a acomodação geométrica a um solucionador simbólico formal.
    \item \textbf{Ponto de Ruptura:} O método foca unicamente em cenas 2D decorativas e estáticas, assemelhando-se a dioramas visuais planos. Ele desconsidera por completo as exigências do \textit{Level Design} interativo de fases jogáveis: não modela a locomoção física de agentes cinemáticos, não calcula Malhas de Navegação (NavMesh), não suporta ambientes modulares conectados por corredores nem trata a progressão funcional e a solvabilidade lógica de quebra-cabeças mecânicos.
\end{itemize}

% ---------------------------------------------------------------------------
\subsection{Síntese de Cenas 3D via Código Intermediário e Renderização \textit{Offline}}
\label{sec:relacionados_codigo_render}
% ---------------------------------------------------------------------------

Diante da constatação empírica de que os LLMs falham ao emitir diretamente listas de coordenadas métricas $(x, y, z)$, uma linha de pesquisa buscou explorar a síntese de código como linguagem intermediária, combinando geradores neurais a suítes de modelagem e renderização tridimensional \textit{offline}.

O principal expoente desse paradigma é o sistema \textit{SceneCraft}, desenvolvido por Hu et al. \cite{hu2024scenecraft} e publicado na ICML 2024.
\begin{itemize}
    \item \textbf{Mecanismo:} O \textit{SceneCraft} estrutura um agente autônomo baseado em LLM encarregado de traduzir descrições em linguagem natural em código executável em linguagem Python, o qual invoca comandos da API interna do software de computação gráfica Blender. O código programa a instanciação de primitivas e malhas 3D, ajustando transformações no espaço contínuo. Para mitigar colisões visuais e erros de proporção, o sistema fecha um laço de retroalimentação multimodal (\textit{feedback loop}): uma câmera virtual renderiza a cena em imagens bidimensionais a partir de múltiplos ângulos, as quais são submetidas a um Modelo de Visão e Linguagem (VLM); o VLM inspeciona os artefatos visuais e devolve críticas em texto para que o LLM reescreva e refine o código Python iterativamente.
    \item \textbf{Mérito:} O emprego de código como camada intermediária formaliza as operações de posicionamento, contornando a emissão de números puros sem contexto sintático. A incorporação do feedback visual via VLM introduz uma verificação empírica automatizada de sobreposições visíveis na renderização.
    \item \textbf{Ponto de Ruptura:} O sistema opera sob um fluxo estritamente \textit{offline} e computacionalmente oneroso, consumindo minutos de processamento por cena devido à renderização fotorrealista e às chamadas de modelos multimodais de alta escala, o que inviabiliza sua aplicação em tempo de execução ou como ferramenta interativa em motores de jogos. Fundamentalmente, o artefato resultante do \textit{SceneCraft} é uma cena gráfica estática e contemplativa: não há Malha de Navegação (NavMesh), não há conectores ou portas funcionais, inexiste física de colisão de tempo real com agentes jogáveis e não há suporte à progressão de objetivos ou validação de solvabilidade de fases.
\end{itemize}

Mais recentemente, Yang et al. \cite{yang2026codeasroomgenerating3drooms} apresentaram o estudo exploratório \textit{Code-as-Room: Generating 3D Rooms from Top-Down View Images via Agentic Code Synthesis}.
\begin{itemize}
    \item \textbf{Mecanismo:} O arcabouço recebe como entrada imagens de plantas baixas residenciais (\textit{top-down view}) e coordena agentes de síntese de código especializados em decompor o perímetro do recinto e alocar mobílias e adereços no espaço tridimensional interno por meio de scripts procedurais.
    \item \textbf{Mérito:} Apresenta alta eficácia na acomodação espacial de móveis respeitando os limites perimetrais de paredes em ambientes residenciais fechados, aproveitando a capacidade de abstração modular do código para evitar sobreposições grosseiras.
    \item \textbf{Ponto de Ruptura:} O escopo do \textit{Code-as-Room} é categoricamente confinado a cômodos residenciais individuais e isolados (\textit{single room layout}). O sistema não possui mecanismos para sintetizar masmorras ou fases completas compostas por múltiplas câmaras articuladas por corredores, bifurcações e transições verticais. Além disso, ignora completamente as variáveis de jogabilidade, fluxo de circulação de jogadores, posicionamento de desafios e encadeamento causal de mecânicas de jogo.
\end{itemize}

% ---------------------------------------------------------------------------
\subsection{Planejamento de \textit{Layouts} 3D Guiado por Otimizadores e Restrições}
\label{sec:relacionados_otimizadores}
% ---------------------------------------------------------------------------

Reconhecendo as falhas geométricas crônicas dos métodos puramente neurais, uma terceira família de pesquisas procurou desacoplar a geração em dois estratos: um planejamento de alto nível orientado por LLM e uma etapa de alocação espacial governada por otimizadores matemáticos ou acervos de dados estruturados.

No trabalho de Sun et al. \cite{sun2024roomplanner}, intitulado \textit{RoomPlanner: Explicit Layout Planner for Easier LLM-Driven 3D Room Generation}, os autores abordaram o arranjo interno de recintos tridimensionais.
\begin{itemize}
    \item \textbf{Mecanismo:} O sistema propõe a divisão explícita em dois estágios sequenciais: primeiramente, o LLM atua como um planejador macroestrutural, gerando a disposição relacional dos objetos sem especificar transformações milimétricas; em seguida, um módulo determinístico de recuperação e posicionamento de ativos (\textit{asset retrieval and placement}) realiza o encaixe físico das malhas dentro do ambiente delimitado.
    \item \textbf{Mérito:} Formaliza a necessidade de segregar a concepção estrutural abstrata da resolução geométrica fina, reconhecendo explicitamente que a imposição de cálculos métricos simultâneos sobrecarrega o LLM e eleva as taxas de alucinação espacial.
    \item \textbf{Ponto de Ruptura:} O modelo é estritamente restrito ao arranjo doméstico de mobília (\textit{indoor furniture placement}). Não contempla as exigências de uma fase de jogo: não existem corredores de transição, conectores modulares, zonas de combate, quebra-cabeças com dependências temporais nem requisitos cinemáticos de navegação transitável para agentes autônomos.
\end{itemize}

Em uma abordagem voltada especificamente para fases de jogos, Xu e Verbrugge \cite{xu2025constraint} apresentaram na FDG 2025 o estudo \textit{Constraint Is All You Need: Optimization-Based 3D Level Generation with LLMs}.
\begin{itemize}
    \item \textbf{Mecanismo:} O usuário fornece uma descrição em linguagem natural do nível desejado. O LLM é instruído a traduzir essa intenção em um conjunto formal de restrições espaciais e topológicas codificadas em uma linguagem de domínio proprietária, denominada GLDL (\textit{Game Level Description Language}). Essas restrições são então mapeadas em uma formulação contínua do Problema de Layout de Instalações (FLP --- \textit{Facility Layout Problem}) e resolvidas por otimizadores numéricos de programação não linear com restrições (tais como IPOPT e SLSQP), que calculam posições contínuas e dimensões de blocos retangulares assegurando a ausência de sobreposição e a adjacência especificada.
    \item \textbf{Mérito:} Estabelece um dos patamares mais rigorosos da literatura recente ao substituir o palpite probabilístico de coordenadas por formulações matemáticas clássicas de otimização, garantindo analiticamente a não sobreposição volumétrica ($V_{\text{overlap}} = 0$) e o cumprimento de adjacências mínimas entre volumes abstratos.
    \item \textbf{Ponto de Ruptura:} A técnica opera como um solucionador numérico puramente matemático e estritamente desacoplado de motores de jogos comerciais (\textit{game engines}). Como consequência desse isolamento: (i) as restrições são calculadas sobre caixas delimitadoras ideais, sem levar em consideração as malhas poligonais detalhadas dos \textit{prefabs} tridimensionais reais; (ii) o sistema não executa o cozimento (\textit{bake}) nem a validação de Malhas de Navegação (NavMesh), ignorando parâmetros cinemáticos de agentes reais (como altura de degraus e raio de passagem); e (iii) a formulação do FLP é puramente geométrica e estática, desconsiderando por completo o espaço de estados da progressão lógica do jogo (isto é, não modela nem verifica se o posicionamento de chaves e portas trancadas gera impasses ou bloqueios intransponíveis na dinâmica da partida).
\end{itemize}

Complementarmente, Xu e Verbrugge \cite{xu2025databasedrivenframework3dlevel} publicaram na AIIDE 2025 o artigo \textit{A Database-Driven Framework for 3D Level Generation with LLMs}.
\begin{itemize}
    \item \textbf{Mecanismo:} Os autores estruturaram um repositório relacional de salas e módulos pré-modelados em 3D, contendo anotações semânticas e pontos de conexão validados. O LLM atua na seleção e composição dessas salas a partir de consultas textuais. Para contornar desconexões geradas pelo modelo durante a montagem do nível, o framework incorpora um algoritmo de pós-processamento baseado em heurísticas de reparo topológico, encarregado de inserir corredores adicionais ou rotacionar salas desalinhadas.
    \item \textbf{Mérito:} O reúso de blocos arquiteturais pré-construídos eleva a fidelidade visual local e garante que as unidades atômicas da masmorra sejam internamente válidas.
    \item \textbf{Ponto de Ruptura:} O sistema padece de alta rigidez, dependendo intrinsecamente de um vocabulário fechado de ativos pré-cadastrados no banco de dados. A necessidade de etapas corretivas \textit{ad hoc} de reparo evidencia a incapacidade do gerador neural em preservar a integridade estrutural da fase de forma autônoma. Ademais, o framework não é integrado nativamente ao fluxo de trabalho de um motor gráfico comercial (como ferramenta ou \textit{plugin} de editor), limitando severamente a interatividade e a ergonomia de autoria assistida para os desenvolvedores.
\end{itemize}

% ---------------------------------------------------------------------------
\subsection{Integração com Motores de Jogos e Arquiteturas Multiagente}
\label{sec:relacionados_motores_multiagente}
% ---------------------------------------------------------------------------

Uma quarta frente de trabalhos tem se concentrado na ponte direta entre LLMs e motores de jogos comerciais (\textit{game engines}), bem como na colaboração entre múltiplos agentes neurais para mitigar erros de geração.

No trabalho de Her et al. \cite{her2025zeroshot3dmapgeneration}, intitulado \textit{Zero-shot 3D Map Generation with LLM Agents: A Dual-Agent Architecture for Procedural Content Generation}, os autores propuseram uma arquitetura multiagente cooperativa.
\begin{itemize}
    \item \textbf{Mecanismo:} O sistema divide o processo de geração entre dois agentes fundamentados em LLMs operando em uma configuração Ator-Crítico (\textit{Actor-Critic}): um agente gerador formula a especificação espacial do mapa a partir do prompt do usuário, e um segundo agente crítico analisa textualmente o rascunho produzido, apontando inconsistências conceituais e devolvendo orientações de refinamento para o gerador em um laço de diálogo \textit{zero-shot}.
    \item \textbf{Mérito:} O trabalho explora a auditoria cruzada sem exigir retreinamento de pesos ou ajustes finos complexos, demonstrando que o diálogo entre agentes pode detectar certas incoerências semânticas e contextuais.
    \item \textbf{Ponto de Ruptura:} A arquitetura incorre no que se denomina na literatura como a \textbf{falácia da auto-validação probabilística}. Ambos os agentes compartilham a mesma natureza conexionista autorregressiva, carecendo intrinsecamente de um modelo físico e geométrico do espaço contínuo. O agente crítico é incapaz de calcular com exatidão se duas caixas delimitadoras tridimensionais colidem no espaço contínuo $\mathbb{R}^3$, não tem acesso a rotinas de colisão física e não pode atestar a conectividade transitável de uma Malha de Navegação poligonal. O ciclo crítico-gerador opera, portanto, em um plano puramente retórico, permanecendo sujeito a alucinações métricas imperceptíveis a ambos os modelos.
\end{itemize}

No âmbito de integração nativa em engines comerciais de alta fidelidade, destaca-se o trabalho de Tang et al. \cite{tang2025unrealllm}, intitulado \textit{UnrealLLM: Towards Highly Controllable and Interactable 3D Scene Generation by LLM-powered Procedural Content Generation}, publicado na ACL 2025.
\begin{itemize}
    \item \textbf{Mecanismo:} O \textit{UnrealLLM} conecta um modelo de linguagem à \textit{Unreal Engine 5} (UE5) por meio de um protocolo baseado em mensagens estruturadas em JSON e ferramentas procedurais baseadas em curvas paramétricas (\textit{splines}). O LLM decompõe instruções em linguagem natural em comandos estruturados para instanciar terrenos, relevos, corpos hídricos, densidade de vegetação e trajetórias de vias, permitindo a autoria de cenas interativas controladas pelo motor gráfico.
    \item \textbf{Mérito:} Representa um dos maiores avanços práticos de acoplamento entre modelos de linguagem e um motor de jogos comercial consolidado na indústria, comprovando que a mediação via esquemas estruturados em JSON viabiliza o controle programático de ambientes gráficos de ponta.
    \item \textbf{Ponto de Ruptura:} O escopo do \textit{UnrealLLM} é quase que exclusivamente direcionado a ambientes abertos (\textit{outdoor}) e à dispersão cosmética e realista de ativos ao longo de superfícies contínuas e splines. O sistema não aborda a geração modular de fases fechadas ou masmorras compostas por salas, corredores e conectores discretos. Fundamentalmente, não há tratamento da progressão mecânica de \textit{Level Design}: o sistema não suporta a formalização de grafos causais acíclicos (DAGs) com pré-requisitos mecânicos (chaves, portas trancadas e interruptores), carecendo de rotinas determinísticas para comprovar a solvabilidade dinâmica da fase no espaço de estados e a integridade de rotas navegáveis em ambientes interiores fechados.
\end{itemize}

Em paralelo, pesquisas complementares como as de Regmi \cite{regmi2025system} e Re \cite{re2026integration} investigaram o emprego de LLMs acoplados a motores de jogos (especialmente Unity). Regmi \cite{regmi2025system} desenvolveu um sistema para geração procedural de atividades de agentes virtuais baseado na interpretação semântica de cenas tridimensionais, ao passo que Re \cite{re2026integration} explorou técnicas de engenharia de \textit{prompt} voltadas à síntese automática de código para agentes inteligentes no desenvolvimento de jogos. Embora ambos os trabalhos atestem a crescente permeabilidade dos LLMs nas rotinas de motores de jogos, seus focos concentram-se estritamente na orquestração de comportamentos de NPCs e na automação de scripts utilitários, sem abordar a síntese procedural da topologia e da geometria de fases tridimensionais, nem os problemas formais de colisão física e solvabilidade mecânica.

% ---------------------------------------------------------------------------
\subsection{Diagnósticos Empíricos de Falhas Espaciais e Percepção do \textit{Level Designer}}
\label{sec:relacionados_diagnosticos}
% ---------------------------------------------------------------------------

Para além das arquiteturas generativas, a consolidação recente da literatura trouxe investigações empíricas e relatos da indústria que evidenciam com rigor metodológico as fragilidades dos geradores conexionistas diretos e a recepção dos profissionais de desenvolvimento.

O estudo empírico de maior impacto nesse domínio foi conduzido por Feng et al. \cite{feng2026repurposing3dgenerativemodel}, sob o título \textit{Repurposing 3D Generative Model for Autoregressive Layout Generation}.
\begin{itemize}
    \item \textbf{Diagnóstico Técnico:} Os autores realizaram testes quantitativos exaustivos sobre a capacidade de modelos generativos neurais e transformadores em emitir diretamente arranjos espaciais tridimensionais. A conclusão central da pesquisa foi formalizada na máxima: \textbf{``Layout as Language: Semantic YES, Physical NO''} (Layout como Linguagem: Sim para Semântica, Não para Física).
    \item \textbf{Evidências Empíricas:} Os experimentos comprovaram que, conquanto os modelos alcancem escores elevados de plausibilidade semântica (como posicionar cadeiras próximas a mesas e associar objetos contextuais corretamente em texto), a geração numérica direta de posições e rotações 3D incorre em falhas geométricas frequentes e inaceitáveis: taxas severas de sobreposição volumétrica não autorizada (\textit{clipping} massivo), elementos suspensos flutuando sem apoio estrutural (\textit{floating objects}), violações perimétricas de colisão contra paredes e desvios de escala que quebram as restrições físicas mínimas do ambiente euclidiano.
    \item \textbf{Implicação para a Área:} Esse diagnóstico confirma teoricamente que a tentativa de delegar o cálculo exato de coordenadas geométricas contínuas a modelos estatísticos autorregressivos constitui uma premissa arquitetural falha, demandando imperativamente que as variáveis métricas sejam governadas por processos determinísticos externos com respaldo nas leis físicas e espaciais.
\end{itemize}

Sob a perspectiva humana e profissional do desenvolvimento de jogos, Bandeali \cite{bandeali2026generative} publicou o relatório investigativo \textit{Generative AI \& The Level Designer: An Exploratory Report of Tools, Workflows, and Challenges}.
\begin{itemize}
    \item \textbf{Diagnóstico da Indústria:} O autor conduziu um estudo de campo detalhado com \textit{level designers} profissionais em estúdios de desenvolvimento de jogos sobre a integração de ferramentas de inteligência artificial generativa em suas rotinas de trabalho.
    \item \textbf{Desafios e Rejeição de Ferramentas Autônomas:} O estudo revelou um ceticismo generalizado da indústria em relação a geradores autônomos do tipo ``caixa-preta'' que entregam cenas fechadas sem controle intermediário. Os profissionais relataram que o tempo despendido para auditar, desmanchar e corrigir as alucinações geométricas, caminhos bloqueados e quebras de ritmo gerados por ferramentas puramente neurais frequentemente supera o tempo necessário para criar a fase manualmente do zero.
    \item \textbf{Requisitos Mandatórios da Indústria:} O relatório conclui que as ferramentas generativas assistidas só alcançam viabilidade e adoção real na indústria quando atendem a três pilares: (i) operação no paradigma de \textbf{iniciativa mista} (\textit{mixed-initiative}), mantendo o designer no controle criativo através de parâmetros claros de edição; (ii) integração nativa e ergonômica ao ambiente do motor de jogo (\textit{game engine}) utilizado pelo estúdio; e (iii) existência de salvaguardas algorítmicas determinísticas que garantam categoricamente que o artefato gerado seja funcional, desprovido de sobreposições volumétricas destrutivas e com navegabilidade atestada.
\end{itemize}

Essa demanda da indústria encontra eco direto em estudos seminais e contemporâneos de iniciativa mista e autoria assistida, tais como os trabalhos de Lai, Latham e Leymarie \cite{lai2020towards} sobre os três pilares de ferramentas amigáveis para a indústria, de Walton, Rahat e Stovold \cite{walton2020evaluating} sobre avaliação cega de ferramentas em estúdios de jogos, e de Wang et al. \cite{wang2026designerlyloop} sobre o alinhamento da intenção do designer com LLMs por meio de circuitos iterativos de raciocínio.

% ---------------------------------------------------------------------------
\subsection{Quadro Comparativo do Estado da Arte}
\label{sec:relacionados_quadro_comparativo}
% ---------------------------------------------------------------------------

Para consolidar as distinções metodológicas e explicitar os limites de cada linha de investigação analisada, a Tabela~\ref{tab:quadro_comparativo_relacionados} (Quadro comparativo do estado da arte) estabelece uma comparação sistemática e multidimensional entre os 11 principais sistemas da literatura e a proposta delineada no presente TCC. A matriz avaliativa cruza os trabalhos frente a sete dimensões técnicas rigorosas:
\begin{enumerate}
    \item \textbf{Paradigma Arquitetural:} a taxonomia do pipeline (puramente neural, multiagente, assistido por otimizador numérico ou híbrido neuro-simbólico);
    \item \textbf{Espaço de Representação:} dimensionalidade do artefato produzido (2D discreto vs. 3D contínuo);
    \item \textbf{Representação Intermediária:} a estrutura de dados empregada entre a entrada e a saída (grade ASCII, triplas relacionais, código imperativo, formulação de restrições ou grafo serializado);
    \item \textbf{Garantia Anti-Clipping:} o mecanismo formal que assegura a não sobreposição física de malhas sólidas no espaço contínuo ($V_{\text{overlap}} = 0$);
    \item \textbf{Navegabilidade Dinâmica:} capacidade de gerar e validar malhas de navegação contínuas (NavMesh) com suporte a agentes cinemáticos sob física da engine;
    \item \textbf{Solvabilidade Mecânica:} capacidade de modelar e comprovar a alcançabilidade dinâmica de quebra-cabeças lógicos (como sistemas de chaves e fechaduras no espaço de estados);
    \item \textbf{Integração com Game Engine:} se o sistema opera incorporado como ferramenta nativa ao editor de um motor de jogos comercial.
\end{enumerate}

\begin{table}[H]
\centering
\caption{Quadro comparativo multidimensional do estado da arte em geração de fases e cenas com LLMs.}
\label{tab:quadro_comparativo_relacionados}
\resizebox{\linewidth}{!}{%
\begin{tabular}{@{}llllllll@{}}
\toprule
\textbf{Trabalho / Sistema} & \textbf{Paradigma} & \textbf{Espaço} & \textbf{Intermediário} & \textbf{Anti-Clipping} & \textbf{Navegabilidade} & \textbf{Solvabilidade} & \textbf{Game Engine} \\ \midrule
Todd et al. \cite{todd2023level} & LLM Puro (ICL) & 2D & Grade ASCII & Não & Não & Baixa (Probabilística) & Não (Simulador) \\
Word2World \cite{nasir2024word2worldgeneratingstoriesworlds} & LLM + Busca & 2D & Tilemap Discreto & Não (Grid) & $A^*$ em Grid 2D & Não & Não (Python) \\
Narrative-to-Scene \cite{chen2025narrative} & LLM + CSP & 2D & Triplas Semânticas & Heurística (2D) & Não & Não se aplica & Não \\ \midrule
SceneCraft \cite{hu2024scenecraft} & LLM + VLM & 3D & Código Python & Feedback Visual & Não & Não se aplica & Não (Blender Offline) \\
Code-as-Room \cite{yang2026codeasroomgenerating3drooms} & Multiagente Código & 3D & Código Python & Heurística Código & Não & Não se aplica & Não (Standalone) \\ \midrule
RoomPlanner \cite{sun2024roomplanner} & LLM + Slot Placement & 3D & Bounding Boxes & Heurística AABB & Não & Não se aplica & Não \\
Xu e Verbrugge (GLDL) \cite{xu2025constraint} & LLM + Otimizador FLP & 3D & GLDL (Restrições) & Numérica Exata & Não (Sem NavMesh) & Não (Estático) & Não (Solver Python) \\
Xu e Verbrugge (DB) \cite{xu2025databasedrivenframework3dlevel} & LLM + Base de Dados & 3D & Grafo Semântico & Heurística Reparo & Parcial (Heurística) & Não & Não \\ \midrule
Her et al. \cite{her2025zeroshot3dmapgeneration} & Multiagente LLM & 3D & Texto Estruturado & Não (Probabilística) & Não & Não & Não \\
UnrealLLM \cite{tang2025unrealllm} & LLM + Procedural UE5 & 3D & JSON / Splines & Física da Engine & Sim (Terreno Aberto) & Não (Sem Chaves) & Sim (Unreal Engine 5) \\
Regmi / Re \cite{regmi2025system, re2026integration} & LLM + Scripting & 3D & Código / Ações & Física da Engine & Parcial (NPCs) & Não se aplica & Sim (Unity) \\ \midrule
\textbf{Proposta deste TCC} & \textbf{Neuro-Simbólico} & \textbf{3D} & \textbf{Grafo JSON (DAG)} & \textbf{Determinística} & \textbf{NavMesh Bake ($A^*$)} & \textbf{Formal (Espaço de Estados)} & \textbf{Sim (Godot Engine)} \\ \bottomrule
\end{tabular}%
}
\vspace{0.1cm}
\begin{flushleft}
\footnotesize{\textbf{Fonte:} Elaborado pelo autor (2026).}
\end{flushleft}
\end{table}

A inspeção detalhada da Tabela~\ref{tab:quadro_comparativo_relacionados} permite deduzir regularidades estruturais claras sobre a literatura contemporânea:
\begin{itemize}
    \item As abordagens bidimensionais \cite{todd2023level, nasir2024word2worldgeneratingstoriesworlds, chen2025narrative} operam satisfatoriamente em representações discretas e permitem a execução de buscas heurísticas simples em grade, mas são fundamentalmente incapazes de transpor suas regras para a complexidade contínua do espaço métrico tridimensional $\mathbb{R}^3$, ignorando volumes de colisão, conectores angulares e relevo vertical;
    \item Os sistemas baseados em síntese de código para suítes de modelagem \textit{offline} \cite{hu2024scenecraft, yang2026codeasroomgenerating3drooms} atingem qualidade estética relevante, mas operam em tempos de processamento proibitivos e geram modelos estáticos e contemplativos, desprovidos de qualquer atributo de jogabilidade, malhas de navegação ou lógica de masmorras;
    \item As abordagens fundamentadas em solucionadores numéricos \cite{xu2025constraint, xu2025databasedrivenframework3dlevel} oferecem sólidas garantias de não colisão volumétrica em modelos ideais, porém permanecem confinadas a processos analíticos desacoplados de motores gráficos comerciais, falhando em gerar Malhas de Navegação e em validar a solvabilidade mecânica dinâmica em tempo de jogo;
    \item As soluções integradas a motores de jogos \cite{tang2025unrealllm, regmi2025system} aproveitam o potencial de engines consagradas, porém concentram-se na dispersão superficial de folhagens e terrenos em ambientes abertos ou no controle comportamental de NPCs, deixando desassistida a síntese modular de masmorras com progressão condicional chave-porta.
\end{itemize}

% ---------------------------------------------------------------------------
\subsection{Síntese da Lacuna na Literatura e Posicionamento da Proposta}
\label{sec:relacionados_lacuna_posicionamento}
% ---------------------------------------------------------------------------

A análise comparativa e o mapeamento dialético dos trabalhos correlatos evidenciam a existência de uma \textbf{tríplice lacuna científica} na literatura contemporânea de Geração Procedural de Fases assistida por Inteligência Artificial:

\begin{enumerate}
    \item \textbf{A falácia da auto-validação probabilística e a ausência de ancoragem física determinística:} Tentativas de utilizar agentes LLMs para auditar outros LLMs em tarefas espaciais \cite{her2025zeroshot3dmapgeneration} incorrem em uma circularidade conexionista estocástica. Modelos autorregressivos não possuem modelo de mundo físico incorporado nem aparato analítico para solucionar geometria contínua euclidiana. Sem um validador simbólico e determinístico operando sobre estruturas de dados espaciais precisas (como árvores de delimitadores e malhas poligonais), é teoricamente inviável assegurar a ausência de sobreposição física e a coerência métrica de ambientes tridimensionais, conforme comprovado empiricamente por Feng et al. \cite{feng2026repurposing3dgenerativemodel}.
    
    \item \textbf{A dissociação entre estética cosmética e arquitetura funcional de \textit{Level Design}:} O estado da arte encontra-se polarizado entre geradores puramente visuais, focados em renderizações fotorrealistas ou arranjos de mobília estática \cite{hu2024scenecraft, sun2024roomplanner}, e modelos abstratos de grafos que ignoram a transitabilidade real do espaço. Nenhuma proposta existente articula conjuntamente a Malha de Navegação (NavMesh com $A^*$) às dependências dinâmicas temporais de missões (como o problema de chaves e portas trancadas governado por um Grafo Acíclico Dirigido no espaço de estados). Ambientes esteticamente impressionantes tornam-se inúteis para a experiência interativa caso não ofereçam garantia formal de completabilidade funcional e navegabilidade.
    
    \item \textbf{O distanciamento entre solucionadores acadêmicos abstratos e motores de jogos comerciais:} Soluções metodologicamente sofisticadas baseadas em restrições \cite{xu2025constraint} operam em isolamento computacional, desconectadas dos pipelines e ferramentas onde os jogos digitais são efetivamente desenvolvidos e produzidos. A literatura carece de arquiteturas neuro-simbólicas integradas nativamente como ferramentas de autoria (\textit{plugins} ou \textit{addons}) em motores de jogos contemporâneos (como o Godot Engine), capazes de oferecer assistência criativa e validação imediata em tempo de desenvolvimento, atendendo categoricamente às demandas de ergonomia e controle expressas pelos profissionais da indústria \cite{bandeali2026generative}.
\end{enumerate}

Para superar de maneira unificada essa tríplice lacuna, o presente Trabalho de Conclusão de Curso propõe uma \textbf{arquitetura híbrida e neuro-simbólica} para a geração assistida de fases tridimensionais modulares:
\begin{itemize}
    \item O estrato cognitivo e semântico de alto nível é delegado ao Modelo de Linguagem de Grande Porte, que atua estritamente como um \textbf{planejador lógico}, traduzindo o prompt em linguagem natural em um grafo causal de progressão estruturado e serializado em formato JSON (contendo salas funcionais, corredores, conexões e pré-requisitos mecânicos de chaves e portas);
    \item O estrato geométrico, físico e mecânico é delegado integralmente a rotinas algorítmicas determinísticas implementadas de forma nativa no motor de jogos \textbf{Godot Engine 4}, que realizam a instanciação espacial contínua de malhas 3D modulares com prevenção rigorosa de colisões (AABB), o cálculo e cozimento programático da Malha de Navegação poligonal (NavMesh), o teste de transitabilidade com $A^*$ e a validação formal da solvabilidade dinâmica no espaço de estados dinâmicos.
\end{itemize}

A formalização dos modelos matemáticos, dos esquemas de dados, dos algoritmos de resolução geométrica e do pipeline de validação implementado na \textit{engine} constitui o escopo detalhado no capítulo de Metodologia que se segue.
